\section{Basic Description}

We shall describe a standard version of type theory called (intensional) Martin-L\"of Type Theory (MLTT).
Set theoretic foundations usually consist of a number of \textit{axioms}, written in the language of 
first-order logic, which is separately assumed beforehand. The logic's deductive system dictates how 
inferences can be made (how the logic behaves), while the \textit{axioms} only define and describe sets and
how they behave. This forms a basis to build mathematics, using sets to describe mathematical
objects.
However, MLTT combines both of these together into the deductive system itself. The
deductive system consists of the primitive notion of `types', the building blocks
for mathematical objects here, and specifies by itself how they ought to behave.

MLTT consists of \textit{rules} for deriving \textit{judgements}, which are either of the form
$ a : A$ meaning `$a$ is a term/instance of type A', or $a \deq b : A$ meaning `$a$ is judgementally equal to $b$',
while both being instances of $A$.
Judgemental equality is a strict notion of equality which classifies objects which are definitionally
the same within the intrinsic language of the theory.
A separate notion `propositional equality', denoted with $=$, is a weaker but more faithful to the 
original notion of equality of mathematical objects.

The language of MLTT is an extension of dependently-typed $\lambda$-calculus, with variables
$x_1, x_2, \ldots$, primitive constants $c, c', c'', \ldots$, and defined constants $f, f', f''$.
The terms $t, t', t'', \ldots$ are given by the grammar 
\begin{align*}
  t \to x \mid \lambda x.t \mid t(t') \mid  c \mid f
\end{align*}

The logic of MLTT comes from treating propositions as types, wherein having an instance/term of a `proposition as a type'
is interpreted as a witness/evidence/proof of the proposition. Other semantics of Type Theory to be described also
afford further interpretation under this lens.~\cite{martin1975intuitionistic}

\subsection{Universes}

Type Theory `postulates' a chain of cumulative universes à la Russell,
  $\U_0, \U_1, \U_2, \ldots$ as primitive constants
with
\begin{itemize}
  \item $\U_i : \U_{i+1}$ for each $i$ 
  \item If $A : \U_i$, then $A : \U_{i+1}$ 
\end{itemize}
Any type in the theory is an instance of some universe in this chain. Each universe is also a type,
and an instance of universes above it in this heirarchy. But unless necessary, we omit the index from
the universe and write $A : \U$.

\textit{\textbf{Remark.} This chain, while indexed by natural numbers, is not related to the Natural 
Numbers that will be defined in the theory. The indices are come the metatheory.
We shall see, the arithmetic object of Natural Numbers $(\N)$ is richer and needs more
information to be properly defined.}
